%%%%%%%%%%%%%%%%%%%%%%%%%%%%%%%%%%%%%%%%%%%%%%%%%%%%%%%%%%%%
% 10.12 SCHEDULING THEORY
% The Composition of Advanced Mathematics
%%%%%%%%%%%%%%%%%%%%%%%%%%%%%%%%%%%%%%%%%%%%%%%%%%%%%%%%%%%%

\section{Scheduling Theory}
\label{sec:scheduling-theory}

\prerequisite{
Optimization fundamentals, combinatorial optimization, integer
programming, dynamic programming, graph theory, queueing concepts,
probability, and basic algorithmic complexity.
}

\learningobjectives{
By the end of this section, the reader should be able to:
\begin{itemize}
    \item define a scheduling problem;
    \item identify jobs, machines, processing times, release dates, and
          due dates;
    \item define completion time, flow time, lateness, tardiness,
          earliness, and makespan;
    \item distinguish single-machine, parallel-machine, flow-shop,
          job-shop, and open-shop environments;
    \item interpret standard three-field scheduling notation;
    \item distinguish preemptive and nonpreemptive scheduling;
    \item model precedence constraints;
    \item construct precedence graphs;
    \item understand critical paths;
    \item formulate common scheduling objectives;
    \item use shortest-processing-time sequencing;
    \item use earliest-due-date sequencing;
    \item understand Smith's weighted shortest-processing-time rule;
    \item recognize Moore--Hodgson scheduling for minimizing the number
          of tardy jobs;
    \item understand list scheduling;
    \item use longest-processing-time scheduling for parallel machines;
    \item understand approximation guarantees for list scheduling;
    \item apply Johnson's rule to a two-machine flow shop;
    \item formulate disjunctive machine constraints;
    \item formulate scheduling problems as mixed-integer programs;
    \item recognize time-indexed formulations;
    \item understand sequence-dependent setup times;
    \item distinguish regular and nonregular scheduling objectives;
    \item understand release-date scheduling;
    \item understand preemptive scheduling conceptually;
    \item recognize shortest-remaining-processing-time scheduling;
    \item understand weighted completion-time objectives;
    \item understand maximum lateness and total tardiness;
    \item recognize scheduling problems that are computationally hard;
    \item derive lower bounds for makespan;
    \item understand branch-and-bound for scheduling;
    \item understand dynamic-programming approaches conceptually;
    \item recognize local-search and neighborhood methods;
    \item understand critical-path scheduling in project networks;
    \item distinguish CPM and PERT conceptually;
    \item compute basic PERT expected activity times;
    \item understand resource-constrained project scheduling conceptually;
    \item recognize workforce and shift scheduling;
    \item connect scheduling theory with manufacturing, transportation,
          computing, logistics, healthcare, and project management.
\end{itemize}
}

%%%%%%%%%%%%%%%%%%%%%%%%%%%%%%%%%%%%%%%%%%%%%%%%%%%%%%%%%%%%
\subsection{What Is Scheduling?}
\label{subsec:what-is-scheduling}
%%%%%%%%%%%%%%%%%%%%%%%%%%%%%%%%%%%%%%%%%%%%%%%%%%%%%%%%%%%%

Scheduling determines:

\[
\boxed{
\text{when}
}
\]

and

\[
\boxed{
\text{where}
}
\]

jobs should be processed using limited resources.

A typical scheduling problem contains:

\[
\boxed{
\text{jobs}
+
\text{machines}
+
\text{processing requirements}
+
\text{timing constraints}.
}
\]

The objective is usually to improve some measure of:

\[
\boxed{
\text{time},
}
\]

\[
\boxed{
\text{delay},
}
\]

\[
\boxed{
\text{throughput},
}
\]

or

\[
\boxed{
\text{resource utilization}.
}
\]

%%%%%%%%%%%%%%%%%%%%%%%%%%%%%%%%%%%%%%%%%%%%%%%%%%%%%%%%%%%%
\subsection{Jobs}
\label{subsec:scheduling-jobs}
%%%%%%%%%%%%%%%%%%%%%%%%%%%%%%%%%%%%%%%%%%%%%%%%%%%%%%%%%%%%

Let

\[
J
=
\{1,2,\ldots,n\}
\]

be the set of jobs.

A job may represent:

\[
\boxed{
\text{a manufactured part},
}
\]

\[
\boxed{
\text{a computer task},
}
\]

\[
\boxed{
\text{an aircraft operation},
}
\]

\[
\boxed{
\text{a maintenance request},
}
\]

or

\[
\boxed{
\text{a project activity}.
}
\]

%%%%%%%%%%%%%%%%%%%%%%%%%%%%%%%%%%%%%%%%%%%%%%%%%%%%%%%%%%%%
\subsection{Machines and Resources}
\label{subsec:scheduling-machines}
%%%%%%%%%%%%%%%%%%%%%%%%%%%%%%%%%%%%%%%%%%%%%%%%%%%%%%%%%%%%

Let

\[
M
=
\{1,2,\ldots,m\}
\]

be the set of machines or resources.

A machine may represent:

\[
\boxed{
\text{physical equipment},
}
\]

\[
\boxed{
\text{a worker},
}
\]

\[
\boxed{
\text{a runway},
}
\]

\[
\boxed{
\text{a processor},
}
\]

or

\[
\boxed{
\text{a service station}.
}
\]

%%%%%%%%%%%%%%%%%%%%%%%%%%%%%%%%%%%%%%%%%%%%%%%%%%%%%%%%%%%%
\subsection{Processing Time}
\label{subsec:scheduling-processing-time}
%%%%%%%%%%%%%%%%%%%%%%%%%%%%%%%%%%%%%%%%%%%%%%%%%%%%%%%%%%%%

The processing time of job \(j\) is commonly denoted

\[
\boxed{
p_j.
}
\]

If processing time depends on the machine, write

\[
\boxed{
p_{ij}
}
\]

for the time required to process job \(j\) on machine \(i\).

%%%%%%%%%%%%%%%%%%%%%%%%%%%%%%%%%%%%%%%%%%%%%%%%%%%%%%%%%%%%
\subsection{Start and Completion Times}
\label{subsec:start-completion-times}
%%%%%%%%%%%%%%%%%%%%%%%%%%%%%%%%%%%%%%%%%%%%%%%%%%%%%%%%%%%%

Let

\[
S_j
=
\text{start time of job }j,
\]

and

\[
C_j
=
\text{completion time of job }j.
\]

For a nonpreemptive job with processing time \(p_j\),

\[
\boxed{
C_j
=
S_j+p_j.
}
\]

%%%%%%%%%%%%%%%%%%%%%%%%%%%%%%%%%%%%%%%%%%%%%%%%%%%%%%%%%%%%
\subsection{Release Dates}
\label{subsec:scheduling-release-dates}
%%%%%%%%%%%%%%%%%%%%%%%%%%%%%%%%%%%%%%%%%%%%%%%%%%%%%%%%%%%%

A release date

\[
\boxed{
r_j
}
\]

is the earliest time at which job \(j\) may begin.

Thus,

\[
\boxed{
S_j
\geq
r_j.
}
\]

A job cannot be processed before the required material, information, or
resource becomes available.

%%%%%%%%%%%%%%%%%%%%%%%%%%%%%%%%%%%%%%%%%%%%%%%%%%%%%%%%%%%%
\subsection{Due Dates}
\label{subsec:scheduling-due-dates}
%%%%%%%%%%%%%%%%%%%%%%%%%%%%%%%%%%%%%%%%%%%%%%%%%%%%%%%%%%%%

A due date

\[
\boxed{
d_j
}
\]

is the desired completion time of job \(j\).

A job completing after \(d_j\) is late in the ordinary scheduling sense.

Due dates create several important performance measures.

%%%%%%%%%%%%%%%%%%%%%%%%%%%%%%%%%%%%%%%%%%%%%%%%%%%%%%%%%%%%
\subsection{Completion Time}
\label{subsec:scheduling-completion-time}
%%%%%%%%%%%%%%%%%%%%%%%%%%%%%%%%%%%%%%%%%%%%%%%%%%%%%%%%%%%%

The completion time of job \(j\) is simply

\[
\boxed{
C_j.
}
\]

The sum of completion times,

\[
\boxed{
\sum_{j=1}^{n}
C_j,
}
\]

is a common measure of total system congestion when all jobs are available
at time zero.

%%%%%%%%%%%%%%%%%%%%%%%%%%%%%%%%%%%%%%%%%%%%%%%%%%%%%%%%%%%%
\subsection{Flow Time}
\label{subsec:scheduling-flow-time}
%%%%%%%%%%%%%%%%%%%%%%%%%%%%%%%%%%%%%%%%%%%%%%%%%%%%%%%%%%%%

The flow time of job \(j\) is

\[
\boxed{
F_j
=
C_j-r_j.
}
\]

It measures how long the job remains in the system after becoming
available.

If every job has

\[
r_j=0,
\]

then

\[
\boxed{
F_j=C_j.
}
\]

%%%%%%%%%%%%%%%%%%%%%%%%%%%%%%%%%%%%%%%%%%%%%%%%%%%%%%%%%%%%
\subsection{Lateness}
\label{subsec:scheduling-lateness}
%%%%%%%%%%%%%%%%%%%%%%%%%%%%%%%%%%%%%%%%%%%%%%%%%%%%%%%%%%%%

The lateness of job \(j\) is

\[
\boxed{
L_j
=
C_j-d_j.
}
\]

If

\[
L_j>0,
\]

the job finishes late.

If

\[
L_j<0,
\]

the job finishes early.

Thus lateness can be negative.

%%%%%%%%%%%%%%%%%%%%%%%%%%%%%%%%%%%%%%%%%%%%%%%%%%%%%%%%%%%%
\subsection{Tardiness}
\label{subsec:scheduling-tardiness}
%%%%%%%%%%%%%%%%%%%%%%%%%%%%%%%%%%%%%%%%%%%%%%%%%%%%%%%%%%%%

The tardiness of job \(j\) is

\[
\boxed{
T_j
=
\max
\{
C_j-d_j,
0
\}.
}
\]

Unlike lateness,

\[
\boxed{
T_j\geq0.
}
\]

Early completion does not produce negative tardiness.

%%%%%%%%%%%%%%%%%%%%%%%%%%%%%%%%%%%%%%%%%%%%%%%%%%%%%%%%%%%%
\subsection{Earliness}
\label{subsec:scheduling-earliness}
%%%%%%%%%%%%%%%%%%%%%%%%%%%%%%%%%%%%%%%%%%%%%%%%%%%%%%%%%%%%

The earliness of job \(j\) is

\[
\boxed{
E_j
=
\max
\{
d_j-C_j,
0
\}.
}
\]

A job cannot simultaneously be both early and tardy:

\[
\boxed{
E_jT_j=0.
}
\]

%%%%%%%%%%%%%%%%%%%%%%%%%%%%%%%%%%%%%%%%%%%%%%%%%%%%%%%%%%%%
\subsection{Tardy-Job Indicator}
\label{subsec:tardy-job-indicator}
%%%%%%%%%%%%%%%%%%%%%%%%%%%%%%%%%%%%%%%%%%%%%%%%%%%%%%%%%%%%

Define

\[
\boxed{
U_j
=
\begin{cases}
1,
&
C_j>d_j,
\\
0,
&
C_j\leq d_j.
\end{cases}
}
\]

Then

\[
\boxed{
\sum_jU_j
}
\]

is the number of tardy jobs.

%%%%%%%%%%%%%%%%%%%%%%%%%%%%%%%%%%%%%%%%%%%%%%%%%%%%%%%%%%%%
\subsection{Makespan}
\label{subsec:scheduling-makespan}
%%%%%%%%%%%%%%%%%%%%%%%%%%%%%%%%%%%%%%%%%%%%%%%%%%%%%%%%%%%%

The makespan is the completion time of the final job:

\[
\boxed{
C_{\max}
=
\max_j
C_j.
}
\]

Minimizing makespan seeks to finish all work as early as possible.

This objective is especially important for parallel-machine and shop
scheduling.

%%%%%%%%%%%%%%%%%%%%%%%%%%%%%%%%%%%%%%%%%%%%%%%%%%%%%%%%%%%%
\subsection{Maximum Lateness}
\label{subsec:maximum-lateness}
%%%%%%%%%%%%%%%%%%%%%%%%%%%%%%%%%%%%%%%%%%%%%%%%%%%%%%%%%%%%

Maximum lateness is

\[
\boxed{
L_{\max}
=
\max_j
L_j.
}
\]

Minimizing \(L_{\max}\) protects against the worst due-date deviation.

%%%%%%%%%%%%%%%%%%%%%%%%%%%%%%%%%%%%%%%%%%%%%%%%%%%%%%%%%%%%
\subsection{Total Tardiness}
\label{subsec:total-tardiness}
%%%%%%%%%%%%%%%%%%%%%%%%%%%%%%%%%%%%%%%%%%%%%%%%%%%%%%%%%%%%

Total tardiness is

\[
\boxed{
\sum_{j=1}^{n}
T_j.
}
\]

A weighted version is

\[
\boxed{
\sum_{j=1}^{n}
w_jT_j,
}
\]

where \(w_j\) represents the importance of job \(j\).

%%%%%%%%%%%%%%%%%%%%%%%%%%%%%%%%%%%%%%%%%%%%%%%%%%%%%%%%%%%%
\subsection{Weighted Completion Time}
\label{subsec:weighted-completion-time}
%%%%%%%%%%%%%%%%%%%%%%%%%%%%%%%%%%%%%%%%%%%%%%%%%%%%%%%%%%%%

If jobs have weights

\[
w_j>0,
\]

the total weighted completion time is

\[
\boxed{
\sum_{j=1}^{n}
w_jC_j.
}
\]

This penalizes delay of high-priority jobs more strongly.

%%%%%%%%%%%%%%%%%%%%%%%%%%%%%%%%%%%%%%%%%%%%%%%%%%%%%%%%%%%%
\subsection{Regular Objectives}
\label{subsec:regular-scheduling-objectives}
%%%%%%%%%%%%%%%%%%%%%%%%%%%%%%%%%%%%%%%%%%%%%%%%%%%%%%%%%%%%

A scheduling objective is regular if delaying any job cannot improve the
objective while all other completion times remain fixed.

Examples include:

\[
\boxed{
C_{\max},
}
\]

\[
\boxed{
\sum C_j,
}
\]

\[
\boxed{
L_{\max},
}
\]

and

\[
\boxed{
\sum T_j.
}
\]

Regular objectives generally provide no incentive to insert unnecessary
idle time when a job is available and feasible.

%%%%%%%%%%%%%%%%%%%%%%%%%%%%%%%%%%%%%%%%%%%%%%%%%%%%%%%%%%%%
\subsection{Nonregular Objectives}
\label{subsec:nonregular-scheduling-objectives}
%%%%%%%%%%%%%%%%%%%%%%%%%%%%%%%%%%%%%%%%%%%%%%%%%%%%%%%%%%%%

Objectives involving earliness can be nonregular.

For example,

\[
\boxed{
\sum_j
(
\alpha_jE_j+\beta_jT_j
)
}
\]

may penalize completing jobs too early as well as too late.

In such problems, intentional idle time may be optimal.

%%%%%%%%%%%%%%%%%%%%%%%%%%%%%%%%%%%%%%%%%%%%%%%%%%%%%%%%%%%%
\subsection{Preemption}
\label{subsec:scheduling-preemption}
%%%%%%%%%%%%%%%%%%%%%%%%%%%%%%%%%%%%%%%%%%%%%%%%%%%%%%%%%%%%

A preemptive schedule allows a job to be interrupted and resumed later.

A nonpreemptive schedule requires a started operation to continue until
completion.

Preemption is possible in some computer systems but impossible or costly
in many manufacturing processes.

%%%%%%%%%%%%%%%%%%%%%%%%%%%%%%%%%%%%%%%%%%%%%%%%%%%%%%%%%%%%
\subsection{Machine Environments}
\label{subsec:scheduling-machine-environments}
%%%%%%%%%%%%%%%%%%%%%%%%%%%%%%%%%%%%%%%%%%%%%%%%%%%%%%%%%%%%

Important scheduling environments include:

\[
\boxed{
\text{single machine},
}
\]

\[
\boxed{
\text{parallel machines},
}
\]

\[
\boxed{
\text{flow shop},
}
\]

\[
\boxed{
\text{job shop},
}
\]

and

\[
\boxed{
\text{open shop}.
}
\]

Each environment imposes different technological restrictions.

%%%%%%%%%%%%%%%%%%%%%%%%%%%%%%%%%%%%%%%%%%%%%%%%%%%%%%%%%%%%
\subsection{Single-Machine Scheduling}
\label{subsec:single-machine-scheduling}
%%%%%%%%%%%%%%%%%%%%%%%%%%%%%%%%%%%%%%%%%%%%%%%%%%%%%%%%%%%%

In single-machine scheduling,

\[
\boxed{
m=1.
}
\]

At most one job can be processed at a time.

The main decision is often the sequence in which jobs should be processed.

%%%%%%%%%%%%%%%%%%%%%%%%%%%%%%%%%%%%%%%%%%%%%%%%%%%%%%%%%%%%
\subsection{Identical Parallel Machines}
\label{subsec:identical-parallel-machines}
%%%%%%%%%%%%%%%%%%%%%%%%%%%%%%%%%%%%%%%%%%%%%%%%%%%%%%%%%%%%

For identical parallel machines, each job has the same processing time on
every machine.

This environment is commonly denoted

\[
\boxed{
P_m.
}
\]

The scheduling decision determines both:

\[
\boxed{
\text{machine assignment}
}
\]

and

\[
\boxed{
\text{job sequence}.
}
\]

%%%%%%%%%%%%%%%%%%%%%%%%%%%%%%%%%%%%%%%%%%%%%%%%%%%%%%%%%%%%
\subsection{Uniform Parallel Machines}
\label{subsec:uniform-parallel-machines}
%%%%%%%%%%%%%%%%%%%%%%%%%%%%%%%%%%%%%%%%%%%%%%%%%%%%%%%%%%%%

Uniform machines have different speeds.

If machine \(i\) has speed

\[
s_i,
\]

then a job with workload \(p_j\) may require

\[
\boxed{
\frac{p_j}{s_i}
}
\]

time on machine \(i\).

This environment is commonly denoted

\[
\boxed{
Q_m.
}
\]

%%%%%%%%%%%%%%%%%%%%%%%%%%%%%%%%%%%%%%%%%%%%%%%%%%%%%%%%%%%%
\subsection{Unrelated Parallel Machines}
\label{subsec:unrelated-parallel-machines}
%%%%%%%%%%%%%%%%%%%%%%%%%%%%%%%%%%%%%%%%%%%%%%%%%%%%%%%%%%%%

For unrelated machines, processing time may depend arbitrarily on both the
job and the machine:

\[
\boxed{
p_{ij}.
}
\]

This environment is commonly denoted

\[
\boxed{
R_m.
}
\]

It is substantially more general than identical-machine scheduling.

%%%%%%%%%%%%%%%%%%%%%%%%%%%%%%%%%%%%%%%%%%%%%%%%%%%%%%%%%%%%
\subsection{Flow-Shop Scheduling}
\label{subsec:flow-shop-scheduling}
%%%%%%%%%%%%%%%%%%%%%%%%%%%%%%%%%%%%%%%%%%%%%%%%%%%%%%%%%%%%

In a flow shop, every job visits machines in the same order.

For example,

\[
\boxed{
M_1
\to
M_2
\to
M_3.
}
\]

A job cannot begin its operation on \(M_2\) until it completes its
operation on \(M_1\).

%%%%%%%%%%%%%%%%%%%%%%%%%%%%%%%%%%%%%%%%%%%%%%%%%%%%%%%%%%%%
\subsection{Job-Shop Scheduling}
\label{subsec:job-shop-scheduling}
%%%%%%%%%%%%%%%%%%%%%%%%%%%%%%%%%%%%%%%%%%%%%%%%%%%%%%%%%%%%

In a job shop, different jobs may require different machine routes.

For example,

\[
J_1:
M_1
\to
M_3
\to
M_2,
\]

while

\[
J_2:
M_2
\to
M_1.
\]

This flexibility makes job-shop scheduling computationally difficult.

%%%%%%%%%%%%%%%%%%%%%%%%%%%%%%%%%%%%%%%%%%%%%%%%%%%%%%%%%%%%
\subsection{Open-Shop Scheduling}
\label{subsec:open-shop-scheduling}
%%%%%%%%%%%%%%%%%%%%%%%%%%%%%%%%%%%%%%%%%%%%%%%%%%%%%%%%%%%%

In an open shop, each job may require processing on several machines, but
the order of those operations is not fixed in advance.

The schedule determines both:

\[
\boxed{
\text{operation ordering}
}
\]

and

\[
\boxed{
\text{resource timing}.
}
\]

%%%%%%%%%%%%%%%%%%%%%%%%%%%%%%%%%%%%%%%%%%%%%%%%%%%%%%%%%%%%
\subsection{Three-Field Scheduling Notation}
\label{subsec:three-field-scheduling-notation}
%%%%%%%%%%%%%%%%%%%%%%%%%%%%%%%%%%%%%%%%%%%%%%%%%%%%%%%%%%%%

A scheduling problem is often represented using

\[
\boxed{
\alpha
\mid
\beta
\mid
\gamma.
}
\]

Here:

\[
\alpha
=
\text{machine environment},
\]

\[
\beta
=
\text{processing restrictions},
\]

and

\[
\gamma
=
\text{objective}.
\]

This is commonly called Graham's three-field notation.

%%%%%%%%%%%%%%%%%%%%%%%%%%%%%%%%%%%%%%%%%%%%%%%%%%%%%%%%%%%%
\subsection{Examples of Scheduling Notation}
\label{subsec:scheduling-notation-examples}
%%%%%%%%%%%%%%%%%%%%%%%%%%%%%%%%%%%%%%%%%%%%%%%%%%%%%%%%%%%%

For example,

\[
\boxed{
1
\mid
\mid
\sum C_j
}
\]

means:

\[
\boxed{
\text{one machine, no additional restrictions, minimize total completion
time}.
}
\]

Similarly,

\[
\boxed{
P_m
\mid
\mid
C_{\max}
}
\]

means:

\[
\boxed{
\text{identical parallel machines, minimize makespan}.
}
\]

%%%%%%%%%%%%%%%%%%%%%%%%%%%%%%%%%%%%%%%%%%%%%%%%%%%%%%%%%%%%
\subsection{Common Restriction Symbols}
\label{subsec:scheduling-restriction-symbols}
%%%%%%%%%%%%%%%%%%%%%%%%%%%%%%%%%%%%%%%%%%%%%%%%%%%%%%%%%%%%

Examples appearing in the middle field include:

\[
\boxed{
r_j
}
\]

for release dates,

\[
\boxed{
\mathrm{prec}
}
\]

for precedence constraints,

\[
\boxed{
\mathrm{pmtn}
}
\]

for preemption,

and

\[
\boxed{
s_{ij}
}
\]

for sequence-dependent setup times.

%%%%%%%%%%%%%%%%%%%%%%%%%%%%%%%%%%%%%%%%%%%%%%%%%%%%%%%%%%%%
\subsection{Shortest Processing Time Rule}
\label{subsec:spt-rule}
%%%%%%%%%%%%%%%%%%%%%%%%%%%%%%%%%%%%%%%%%%%%%%%%%%%%%%%%%%%%

The shortest processing time rule, abbreviated SPT, orders jobs by
nondecreasing processing time:

\[
\boxed{
p_1
\leq
p_2
\leq
\cdots
\leq
p_n.
}
\]

For

\[
\boxed{
1\mid\mid\sum C_j,
}
\]

SPT is optimal.

%%%%%%%%%%%%%%%%%%%%%%%%%%%%%%%%%%%%%%%%%%%%%%%%%%%%%%%%%%%%
\subsection{Why SPT Minimizes Total Completion Time}
\label{subsec:spt-proof}
%%%%%%%%%%%%%%%%%%%%%%%%%%%%%%%%%%%%%%%%%%%%%%%%%%%%%%%%%%%%

Suppose two consecutive jobs \(i\) and \(j\) satisfy

\[
p_i>p_j,
\]

but \(i\) is scheduled before \(j\).

Ignoring the completion time accumulated before them, their contribution
to total completion time is

\[
p_i
+
(p_i+p_j)
=
2p_i+p_j.
\]

If their order is reversed, the contribution is

\[
p_j
+
(p_j+p_i)
=
p_i+2p_j.
\]

The difference is

\[
(2p_i+p_j)
-
(p_i+2p_j)
=
p_i-p_j
>
0.
\]

Therefore placing the shorter job first improves the objective.

Repeated pairwise exchanges yield SPT order.

%%%%%%%%%%%%%%%%%%%%%%%%%%%%%%%%%%%%%%%%%%%%%%%%%%%%%%%%%%%%
\subsection{Worked Example: SPT}
\label{subsec:worked-spt}
%%%%%%%%%%%%%%%%%%%%%%%%%%%%%%%%%%%%%%%%%%%%%%%%%%%%%%%%%%%%

\begin{workedexamplebox}{Minimizing Total Completion Time}

Suppose four jobs have processing times

\[
p_A=6,
\qquad
p_B=2,
\qquad
p_C=5,
\qquad
p_D=1.
\]

SPT order is

\[
D,B,C,A.
\]

Completion times are

\[
C_D=1,
\]

\[
C_B=3,
\]

\[
C_C=8,
\]

and

\[
C_A=14.
\]

Therefore,

\[
\sum C_j
=
1+3+8+14
=
26.
\]

\begin{answerbox}
\[
\boxed{
D\to B\to C\to A
}
\]

minimizes total completion time.
\end{answerbox}

\end{workedexamplebox}

%%%%%%%%%%%%%%%%%%%%%%%%%%%%%%%%%%%%%%%%%%%%%%%%%%%%%%%%%%%%
\subsection{Weighted Shortest Processing Time}
\label{subsec:wspt-rule}
%%%%%%%%%%%%%%%%%%%%%%%%%%%%%%%%%%%%%%%%%%%%%%%%%%%%%%%%%%%%

For

\[
\boxed{
1\mid\mid\sum_jw_jC_j,
}
\]

the optimal sequence is obtained by ordering jobs according to

\[
\boxed{
\frac{p_j}{w_j}
}
\]

in nondecreasing order.

Equivalently, order by

\[
\boxed{
\frac{w_j}{p_j}
}
\]

in nonincreasing order.

This is often called Smith's rule or the weighted shortest processing time
rule.

%%%%%%%%%%%%%%%%%%%%%%%%%%%%%%%%%%%%%%%%%%%%%%%%%%%%%%%%%%%%
\subsection{Pairwise Interchange for Smith's Rule}
\label{subsec:smith-rule-proof}
%%%%%%%%%%%%%%%%%%%%%%%%%%%%%%%%%%%%%%%%%%%%%%%%%%%%%%%%%%%%

Suppose jobs \(i\) and \(j\) are adjacent.

Scheduling \(i\) before \(j\) contributes, apart from prior completion
time,

\[
w_i p_i
+
w_j(p_i+p_j).
\]

Scheduling \(j\) before \(i\) contributes

\[
w_j p_j
+
w_i(p_j+p_i).
\]

Placing \(i\) before \(j\) is no worse when

\[
w_jp_i
\leq
w_ip_j.
\]

Equivalently,

\[
\boxed{
\frac{p_i}{w_i}
\leq
\frac{p_j}{w_j}.
}
\]

%%%%%%%%%%%%%%%%%%%%%%%%%%%%%%%%%%%%%%%%%%%%%%%%%%%%%%%%%%%%
\subsection{Earliest Due Date Rule}
\label{subsec:edd-rule}
%%%%%%%%%%%%%%%%%%%%%%%%%%%%%%%%%%%%%%%%%%%%%%%%%%%%%%%%%%%%

The earliest due date rule, abbreviated EDD, orders jobs by nondecreasing
due date:

\[
\boxed{
d_1
\leq
d_2
\leq
\cdots
\leq
d_n.
}
\]

For

\[
\boxed{
1\mid\mid L_{\max},
}
\]

EDD minimizes maximum lateness.

%%%%%%%%%%%%%%%%%%%%%%%%%%%%%%%%%%%%%%%%%%%%%%%%%%%%%%%%%%%%
\subsection{Worked Example: EDD}
\label{subsec:worked-edd}
%%%%%%%%%%%%%%%%%%%%%%%%%%%%%%%%%%%%%%%%%%%%%%%%%%%%%%%%%%%%

\begin{workedexamplebox}{Minimizing Maximum Lateness}

Suppose jobs have

\[
(p_A,d_A)=(4,10),
\]

\[
(p_B,d_B)=(3,5),
\]

\[
(p_C,d_C)=(2,8).
\]

EDD order is

\[
B,C,A.
\]

Completion times are

\[
C_B=3,
\]

\[
C_C=5,
\]

\[
C_A=9.
\]

Lateness values are

\[
L_B=3-5=-2,
\]

\[
L_C=5-8=-3,
\]

\[
L_A=9-10=-1.
\]

Therefore,

\[
\boxed{
L_{\max}=-1.
}
\]

\begin{answerbox}
\[
\boxed{
B\to C\to A
}
\]

is an EDD schedule.
\end{answerbox}

\end{workedexamplebox}

%%%%%%%%%%%%%%%%%%%%%%%%%%%%%%%%%%%%%%%%%%%%%%%%%%%%%%%%%%%%
\subsection{Moore--Hodgson Algorithm}
\label{subsec:moore-hodgson}
%%%%%%%%%%%%%%%%%%%%%%%%%%%%%%%%%%%%%%%%%%%%%%%%%%%%%%%%%%%%

For

\[
\boxed{
1\mid\mid\sum U_j,
}
\]

the Moore--Hodgson algorithm minimizes the number of tardy jobs.

The algorithm:

\begin{enumerate}
    \item orders jobs by EDD;
    \item processes jobs in that order;
    \item whenever the current sequence becomes tardy, removes the job
          with the largest processing time among jobs currently scheduled;
    \item places removed jobs after the on-time jobs.
\end{enumerate}

%%%%%%%%%%%%%%%%%%%%%%%%%%%%%%%%%%%%%%%%%%%%%%%%%%%%%%%%%%%%
\subsection{Interpretation of Moore--Hodgson}
\label{subsec:moore-hodgson-interpretation}
%%%%%%%%%%%%%%%%%%%%%%%%%%%%%%%%%%%%%%%%%%%%%%%%%%%%%%%%%%%%

When an EDD prefix cannot meet all due dates, removing the longest job
frees the largest amount of processing time.

This maximizes the opportunity for the remaining jobs to finish on time.

%%%%%%%%%%%%%%%%%%%%%%%%%%%%%%%%%%%%%%%%%%%%%%%%%%%%%%%%%%%%
\subsection{Total Tardiness Is Harder}
\label{subsec:total-tardiness-hardness}
%%%%%%%%%%%%%%%%%%%%%%%%%%%%%%%%%%%%%%%%%%%%%%%%%%%%%%%%%%%%

Although EDD optimizes

\[
L_{\max},
\]

it does not generally minimize

\[
\sum_jT_j.
\]

Similarly, Moore--Hodgson minimizes the number of tardy jobs but not their
total tardiness.

The single-machine total-tardiness problem is substantially more
difficult.

%%%%%%%%%%%%%%%%%%%%%%%%%%%%%%%%%%%%%%%%%%%%%%%%%%%%%%%%%%%%
\subsection{Critical Ratio Dispatching Rule}
\label{subsec:critical-ratio-rule}
%%%%%%%%%%%%%%%%%%%%%%%%%%%%%%%%%%%%%%%%%%%%%%%%%%%%%%%%%%%%

A common real-time dispatching priority index is

\[
\boxed{
CR_j
=
\frac{
d_j-t
}{
p_j^{\mathrm{remaining}}
},
}
\]

where \(t\) is the current time.

A small critical ratio indicates greater urgency.

This is a heuristic dispatching rule, not a general optimality theorem.

%%%%%%%%%%%%%%%%%%%%%%%%%%%%%%%%%%%%%%%%%%%%%%%%%%%%%%%%%%%%
\subsection{Other Dispatching Rules}
\label{subsec:other-dispatching-rules}
%%%%%%%%%%%%%%%%%%%%%%%%%%%%%%%%%%%%%%%%%%%%%%%%%%%%%%%%%%%%

Common shop-floor rules include:

\[
\boxed{
\text{FIFO},
}
\]

\[
\boxed{
\text{SPT},
}
\]

\[
\boxed{
\text{EDD},
}
\]

\[
\boxed{
\text{least slack},
}
\]

and

\[
\boxed{
\text{critical ratio}.
}
\]

Different rules optimize different operational priorities.

%%%%%%%%%%%%%%%%%%%%%%%%%%%%%%%%%%%%%%%%%%%%%%%%%%%%%%%%%%%%
\subsection{Slack}
\label{subsec:scheduling-slack}
%%%%%%%%%%%%%%%%%%%%%%%%%%%%%%%%%%%%%%%%%%%%%%%%%%%%%%%%%%%%

At time \(t\), the slack of job \(j\) is

\[
\boxed{
SL_j
=
d_j
-
t
-
p_j^{\mathrm{remaining}}.
}
\]

A negative slack means the job cannot meet its due date even if processed
continuously from the current time.

%%%%%%%%%%%%%%%%%%%%%%%%%%%%%%%%%%%%%%%%%%%%%%%%%%%%%%%%%%%%
\subsection{Precedence Constraints}
\label{subsec:scheduling-precedence}
%%%%%%%%%%%%%%%%%%%%%%%%%%%%%%%%%%%%%%%%%%%%%%%%%%%%%%%%%%%%

Suppose job \(i\) must finish before job \(j\) begins.

Write

\[
\boxed{
i\prec j.
}
\]

Then

\[
\boxed{
S_j
\geq
C_i.
}
\]

For nonpreemptive processing,

\[
\boxed{
S_j
\geq
S_i+p_i.
}
\]

%%%%%%%%%%%%%%%%%%%%%%%%%%%%%%%%%%%%%%%%%%%%%%%%%%%%%%%%%%%%
\subsection{Precedence Graph}
\label{subsec:precedence-graph}
%%%%%%%%%%%%%%%%%%%%%%%%%%%%%%%%%%%%%%%%%%%%%%%%%%%%%%%%%%%%

Precedence constraints can be represented by a directed acyclic graph.

A directed edge

\[
i\to j
\]

means

\[
\boxed{
i
\text{ must precede }
j.
}
\]

A feasible schedule must respect a topological ordering of the graph.

%%%%%%%%%%%%%%%%%%%%%%%%%%%%%%%%%%%%%%%%%%%%%%%%%%%%%%%%%%%%
\subsection{Critical Path}
\label{subsec:scheduling-critical-path}
%%%%%%%%%%%%%%%%%%%%%%%%%%%%%%%%%%%%%%%%%%%%%%%%%%%%%%%%%%%%

In a project network, a path from project start to project completion has
length equal to the sum of activity durations along that path.

The critical path is a longest-duration path.

Its length gives a lower bound on project completion time:

\[
\boxed{
C_{\max}
\geq
\text{critical-path length}.
}
\]

%%%%%%%%%%%%%%%%%%%%%%%%%%%%%%%%%%%%%%%%%%%%%%%%%%%%%%%%%%%%
\subsection{Earliest Start Times}
\label{subsec:earliest-start-times}
%%%%%%%%%%%%%%%%%%%%%%%%%%%%%%%%%%%%%%%%%%%%%%%%%%%%%%%%%%%%

For a precedence-constrained project without resource conflicts, the
earliest start time of an activity is determined by the completion of all
its predecessors.

If \(\operatorname{Pred}(j)\) is the predecessor set,

\[
\boxed{
ES_j
=
\max_{i\in\operatorname{Pred}(j)}
EF_i,
}
\]

where

\[
\boxed{
EF_i
=
ES_i+p_i.
}
\]

%%%%%%%%%%%%%%%%%%%%%%%%%%%%%%%%%%%%%%%%%%%%%%%%%%%%%%%%%%%%
\subsection{Latest Start Times}
\label{subsec:latest-start-times}
%%%%%%%%%%%%%%%%%%%%%%%%%%%%%%%%%%%%%%%%%%%%%%%%%%%%%%%%%%%%

Working backward from the project completion time gives latest finish and
latest start times.

For activity \(j\),

\[
\boxed{
LS_j
=
LF_j-p_j.
}
\]

The total slack is

\[
\boxed{
LS_j-ES_j.
}
\]

Critical activities have zero slack.

%%%%%%%%%%%%%%%%%%%%%%%%%%%%%%%%%%%%%%%%%%%%%%%%%%%%%%%%%%%%
\subsection{Critical Path Method}
\label{subsec:critical-path-method}
%%%%%%%%%%%%%%%%%%%%%%%%%%%%%%%%%%%%%%%%%%%%%%%%%%%%%%%%%%%%

The Critical Path Method, abbreviated CPM, assumes deterministic activity
durations.

The basic procedure is:

\begin{enumerate}
    \item construct the precedence network;
    \item perform a forward pass;
    \item determine project completion time;
    \item perform a backward pass;
    \item compute activity slack;
    \item identify the critical path.
\end{enumerate}

%%%%%%%%%%%%%%%%%%%%%%%%%%%%%%%%%%%%%%%%%%%%%%%%%%%%%%%%%%%%
\subsection{PERT}
\label{subsec:pert}
%%%%%%%%%%%%%%%%%%%%%%%%%%%%%%%%%%%%%%%%%%%%%%%%%%%%%%%%%%%%

The Program Evaluation and Review Technique, abbreviated PERT, introduces
uncertain activity durations.

A classical PERT model uses:

\[
a
=
\text{optimistic duration},
\]

\[
m
=
\text{most likely duration},
\]

and

\[
b
=
\text{pessimistic duration}.
\]

%%%%%%%%%%%%%%%%%%%%%%%%%%%%%%%%%%%%%%%%%%%%%%%%%%%%%%%%%%%%
\subsection{PERT Expected Duration}
\label{subsec:pert-expected-duration}
%%%%%%%%%%%%%%%%%%%%%%%%%%%%%%%%%%%%%%%%%%%%%%%%%%%%%%%%%%%%

The classical PERT approximation uses

\[
\boxed{
\mathbb{E}[T]
\approx
\frac{
a+4m+b
}{
6
}.
}
\]

The approximate variance is

\[
\boxed{
\operatorname{Var}(T)
\approx
\left(
\frac{
b-a
}{
6
}
\right)^2.
}
\]

These formulas arise from the traditional beta-distribution approximation
used in PERT.

%%%%%%%%%%%%%%%%%%%%%%%%%%%%%%%%%%%%%%%%%%%%%%%%%%%%%%%%%%%%
\subsection{Worked Example: PERT Duration}
\label{subsec:worked-pert-duration}
%%%%%%%%%%%%%%%%%%%%%%%%%%%%%%%%%%%%%%%%%%%%%%%%%%%%%%%%%%%%

\begin{workedexamplebox}{Expected Activity Duration}

Suppose an activity has:

\[
a=4,
\qquad
m=7,
\qquad
b=10.
\]

Then

\[
\mathbb{E}[T]
\approx
\frac{
4+4(7)+10
}{
6
}.
\]

Thus,

\[
\mathbb{E}[T]
=
\frac{42}{6}
=
7.
\]

The approximate variance is

\[
\left(
\frac{10-4}{6}
\right)^2
=
1.
\]

\begin{answerbox}
\[
\boxed{
\mathbb{E}[T]\approx7,
\qquad
\operatorname{Var}(T)\approx1.
}
\]
\end{answerbox}

\end{workedexamplebox}

%%%%%%%%%%%%%%%%%%%%%%%%%%%%%%%%%%%%%%%%%%%%%%%%%%%%%%%%%%%%
\subsection{Resource Constraints in Projects}
\label{subsec:resource-constrained-projects}
%%%%%%%%%%%%%%%%%%%%%%%%%%%%%%%%%%%%%%%%%%%%%%%%%%%%%%%%%%%%

Critical-path calculations assume that precedence is the main restriction.

In practice, multiple activities may compete for limited resources.

A resource capacity constraint may be written

\[
\boxed{
\sum_j
a_{rj}
y_{jt}
\leq
R_r,
}
\]

where:

\[
a_{rj}
=
\text{units of resource }r\text{ used by job }j,
\]

\[
y_{jt}
=
\text{whether job }j\text{ is active at time }t,
\]

and

\[
R_r
=
\text{available capacity}.
\]

%%%%%%%%%%%%%%%%%%%%%%%%%%%%%%%%%%%%%%%%%%%%%%%%%%%%%%%%%%%%
\subsection{Resource-Constrained Project Scheduling}
\label{subsec:rcpsp}
%%%%%%%%%%%%%%%%%%%%%%%%%%%%%%%%%%%%%%%%%%%%%%%%%%%%%%%%%%%%

The Resource-Constrained Project Scheduling Problem, abbreviated RCPSP,
combines:

\[
\boxed{
\text{precedence constraints}
}
\]

with

\[
\boxed{
\text{limited renewable resources}.
}
\]

A common objective is

\[
\boxed{
\min C_{\max}.
}
\]

The RCPSP is computationally difficult in general.

%%%%%%%%%%%%%%%%%%%%%%%%%%%%%%%%%%%%%%%%%%%%%%%%%%%%%%%%%%%%
\subsection{Parallel-Machine Makespan Scheduling}
\label{subsec:parallel-machine-makespan}
%%%%%%%%%%%%%%%%%%%%%%%%%%%%%%%%%%%%%%%%%%%%%%%%%%%%%%%%%%%%

For identical parallel machines,

\[
\boxed{
P_m
\mid
\mid
C_{\max},
}
\]

jobs must be distributed among \(m\) machines.

The load on machine \(i\) is

\[
\boxed{
L_i
=
\sum_{j\text{ assigned to }i}
p_j.
}
\]

Then

\[
\boxed{
C_{\max}
=
\max_i
L_i.
}
\]

%%%%%%%%%%%%%%%%%%%%%%%%%%%%%%%%%%%%%%%%%%%%%%%%%%%%%%%%%%%%
\subsection{Lower Bounds for Parallel-Machine Makespan}
\label{subsec:parallel-machine-lower-bounds}
%%%%%%%%%%%%%%%%%%%%%%%%%%%%%%%%%%%%%%%%%%%%%%%%%%%%%%%%%%%%

Any schedule must satisfy

\[
\boxed{
C_{\max}
\geq
\frac{
\sum_jp_j
}{
m
}.
}
\]

Also,

\[
\boxed{
C_{\max}
\geq
\max_jp_j.
}
\]

Therefore a basic lower bound is

\[
\boxed{
LB
=
\max
\left\{
\frac{\sum_jp_j}{m},
\;
\max_jp_j
\right\}.
}
\]

%%%%%%%%%%%%%%%%%%%%%%%%%%%%%%%%%%%%%%%%%%%%%%%%%%%%%%%%%%%%
\subsection{List Scheduling}
\label{subsec:list-scheduling}
%%%%%%%%%%%%%%%%%%%%%%%%%%%%%%%%%%%%%%%%%%%%%%%%%%%%%%%%%%%%

List scheduling processes jobs in a specified list order.

Whenever a machine becomes available, assign the next unscheduled job to
that machine.

For identical parallel machines and makespan minimization, any list order
produces a schedule satisfying the classical bound

\[
\boxed{
C_{\max}^{LS}
\leq
\left(
2-\frac1m
\right)
C_{\max}^*.
}
\]

%%%%%%%%%%%%%%%%%%%%%%%%%%%%%%%%%%%%%%%%%%%%%%%%%%%%%%%%%%%%
\subsection{Longest Processing Time Rule}
\label{subsec:lpt-rule}
%%%%%%%%%%%%%%%%%%%%%%%%%%%%%%%%%%%%%%%%%%%%%%%%%%%%%%%%%%%%

Longest Processing Time, abbreviated LPT, first sorts jobs so that

\[
\boxed{
p_1
\geq
p_2
\geq
\cdots
\geq
p_n.
}
\]

Then list scheduling is applied.

Large jobs are distributed early, reducing the risk that a large job
creates a poor final imbalance.

%%%%%%%%%%%%%%%%%%%%%%%%%%%%%%%%%%%%%%%%%%%%%%%%%%%%%%%%%%%%
\subsection{LPT Approximation Bound}
\label{subsec:lpt-approximation-bound}
%%%%%%%%%%%%%%%%%%%%%%%%%%%%%%%%%%%%%%%%%%%%%%%%%%%%%%%%%%%%

For

\[
m\geq2,
\]

LPT satisfies the classical bound

\[
\boxed{
C_{\max}^{LPT}
\leq
\left(
\frac43
-
\frac{1}{3m}
\right)
C_{\max}^*.
}
\]

Thus LPT has a stronger worst-case guarantee than arbitrary list
scheduling.

%%%%%%%%%%%%%%%%%%%%%%%%%%%%%%%%%%%%%%%%%%%%%%%%%%%%%%%%%%%%
\subsection{Worked Example: LPT}
\label{subsec:worked-lpt}
%%%%%%%%%%%%%%%%%%%%%%%%%%%%%%%%%%%%%%%%%%%%%%%%%%%%%%%%%%%%

\begin{workedexamplebox}{Three Identical Machines}

Suppose processing times are

\[
8,\ 7,\ 6,\ 5,\ 4,\ 3.
\]

There are

\[
m=3
\]

machines.

LPT assigns:

\[
M_1:
8+3=11,
\]

\[
M_2:
7+4=11,
\]

\[
M_3:
6+5=11.
\]

Therefore,

\begin{answerbox}
\[
\boxed{
C_{\max}=11.
}
\]
\end{answerbox}

Since

\[
\frac{
8+7+6+5+4+3
}{
3
}
=
11,
\]

the lower bound is also \(11\).

Hence the LPT schedule is optimal.

\end{workedexamplebox}

%%%%%%%%%%%%%%%%%%%%%%%%%%%%%%%%%%%%%%%%%%%%%%%%%%%%%%%%%%%%
\subsection{Load Balancing}
\label{subsec:scheduling-load-balancing}
%%%%%%%%%%%%%%%%%%%%%%%%%%%%%%%%%%%%%%%%%%%%%%%%%%%%%%%%%%%%

Parallel-machine makespan scheduling is a load-balancing problem.

Each job contributes workload \(p_j\).

The goal is to partition jobs among machines so that the largest total
load is minimized.

This connects scheduling with partitioning and combinatorial optimization.

%%%%%%%%%%%%%%%%%%%%%%%%%%%%%%%%%%%%%%%%%%%%%%%%%%%%%%%%%%%%
\subsection{Flow-Shop Scheduling}
\label{subsec:flow-shop-details}
%%%%%%%%%%%%%%%%%%%%%%%%%%%%%%%%%%%%%%%%%%%%%%%%%%%%%%%%%%%%

Suppose each job requires processing on two machines in order:

\[
M_1
\to
M_2.
\]

Let

\[
p_{1j}
\]

be processing time of job \(j\) on machine \(1\),

and

\[
p_{2j}
\]

its processing time on machine \(2\).

The problem

\[
\boxed{
F_2
\mid
\mid
C_{\max}
}
\]

has an elegant optimal sequencing rule.

%%%%%%%%%%%%%%%%%%%%%%%%%%%%%%%%%%%%%%%%%%%%%%%%%%%%%%%%%%%%
\subsection{Johnson's Rule}
\label{subsec:johnsons-rule}
%%%%%%%%%%%%%%%%%%%%%%%%%%%%%%%%%%%%%%%%%%%%%%%%%%%%%%%%%%%%

For the two-machine flow-shop makespan problem, Johnson's rule is:

\begin{enumerate}
    \item find the smallest unscheduled processing time among all remaining
          operations;
    \item if it occurs on machine \(1\), place that job in the earliest
          available position;
    \item if it occurs on machine \(2\), place that job in the latest
          available position;
    \item remove the job and repeat.
\end{enumerate}

The resulting sequence minimizes makespan.

%%%%%%%%%%%%%%%%%%%%%%%%%%%%%%%%%%%%%%%%%%%%%%%%%%%%%%%%%%%%
\subsection{Alternative Statement of Johnson's Rule}
\label{subsec:johnsons-rule-partition}
%%%%%%%%%%%%%%%%%%%%%%%%%%%%%%%%%%%%%%%%%%%%%%%%%%%%%%%%%%%%

Partition jobs into:

\[
\boxed{
A
=
\{
j:
p_{1j}\leq p_{2j}
\},
}
\]

and

\[
\boxed{
B
=
\{
j:
p_{1j}>p_{2j}
\}.
}
\]

Schedule jobs in \(A\) first in nondecreasing order of \(p_{1j}\).

Then schedule jobs in \(B\) in nonincreasing order of \(p_{2j}\).

%%%%%%%%%%%%%%%%%%%%%%%%%%%%%%%%%%%%%%%%%%%%%%%%%%%%%%%%%%%%
\subsection{Worked Example: Johnson's Rule}
\label{subsec:worked-johnson}
%%%%%%%%%%%%%%%%%%%%%%%%%%%%%%%%%%%%%%%%%%%%%%%%%%%%%%%%%%%%

\begin{workedexamplebox}{Two-Machine Flow Shop}

Suppose:

\[
\begin{array}{c|cccc}
\text{Job}
&
A
&
B
&
C
&
D
\\
\hline
p_{1j}
&
3
&
8
&
4
&
7
\\
p_{2j}
&
6
&
2
&
5
&
9
\end{array}
\]

Jobs with

\[
p_{1j}\leq p_{2j}
\]

are

\[
A,C,D.
\]

Order these by increasing \(p_{1j}\):

\[
A,C,D.
\]

Job \(B\) belongs to the second group.

Therefore Johnson's sequence is

\begin{answerbox}
\[
\boxed{
A\to C\to D\to B.
}
\]
\end{answerbox}

\end{workedexamplebox}

%%%%%%%%%%%%%%%%%%%%%%%%%%%%%%%%%%%%%%%%%%%%%%%%%%%%%%%%%%%%
\subsection{Permutation Flow Shops}
\label{subsec:permutation-flow-shops}
%%%%%%%%%%%%%%%%%%%%%%%%%%%%%%%%%%%%%%%%%%%%%%%%%%%%%%%%%%%%

In a permutation flow shop, all machines process jobs in the same job
order.

If the sequence is

\[
J_3,J_1,J_2,
\]

then every machine processes jobs in that order.

For more than two machines, makespan minimization is generally much more
difficult than the Johnson-rule case.

%%%%%%%%%%%%%%%%%%%%%%%%%%%%%%%%%%%%%%%%%%%%%%%%%%%%%%%%%%%%
\subsection{Job-Shop Disjunctive Constraints}
\label{subsec:job-shop-disjunctive}
%%%%%%%%%%%%%%%%%%%%%%%%%%%%%%%%%%%%%%%%%%%%%%%%%%%%%%%%%%%%

Suppose two operations \(i\) and \(j\) require the same machine.

They cannot overlap.

Therefore either

\[
\boxed{
S_i+p_i
\leq
S_j,
}
\]

or

\[
\boxed{
S_j+p_j
\leq
S_i.
}
\]

This is a disjunctive constraint.

%%%%%%%%%%%%%%%%%%%%%%%%%%%%%%%%%%%%%%%%%%%%%%%%%%%%%%%%%%%%
\subsection{Binary Sequencing Variable}
\label{subsec:binary-sequencing-variable}
%%%%%%%%%%%%%%%%%%%%%%%%%%%%%%%%%%%%%%%%%%%%%%%%%%%%%%%%%%%%

Define

\[
y_{ij}
=
\begin{cases}
1,
&
i\text{ is scheduled before }j,
\\
0,
&
j\text{ is scheduled before }i.
\end{cases}
\]

A Big-\(M\) formulation is

\[
\boxed{
S_i+p_i
\leq
S_j
+
M(1-y_{ij}),
}
\]

and

\[
\boxed{
S_j+p_j
\leq
S_i
+
My_{ij}.
}
\]

%%%%%%%%%%%%%%%%%%%%%%%%%%%%%%%%%%%%%%%%%%%%%%%%%%%%%%%%%%%%
\subsection{Choosing Big-\(M\) in Scheduling}
\label{subsec:scheduling-big-m}
%%%%%%%%%%%%%%%%%%%%%%%%%%%%%%%%%%%%%%%%%%%%%%%%%%%%%%%%%%%%

The value \(M\) must be large enough to deactivate the appropriate
sequencing inequality.

However, an excessively large value can create:

\[
\boxed{
\text{weak LP relaxations}
}
\]

and

\[
\boxed{
\text{poor numerical conditioning}.
}
\]

A valid scheduling horizon often provides a useful upper bound for \(M\).

%%%%%%%%%%%%%%%%%%%%%%%%%%%%%%%%%%%%%%%%%%%%%%%%%%%%%%%%%%%%
\subsection{Makespan Variable in MILP}
\label{subsec:makespan-variable-milp}
%%%%%%%%%%%%%%%%%%%%%%%%%%%%%%%%%%%%%%%%%%%%%%%%%%%%%%%%%%%%

Introduce variable

\[
C_{\max}.
\]

For every job,

\[
\boxed{
C_j
\leq
C_{\max}.
}
\]

Then minimize

\[
\boxed{
C_{\max}.
}
\]

This converts the maximum completion time into a linear objective.

%%%%%%%%%%%%%%%%%%%%%%%%%%%%%%%%%%%%%%%%%%%%%%%%%%%%%%%%%%%%
\subsection{Tardiness Linearization}
\label{subsec:tardiness-linearization}
%%%%%%%%%%%%%%%%%%%%%%%%%%%%%%%%%%%%%%%%%%%%%%%%%%%%%%%%%%%%

Introduce variable

\[
T_j\geq0.
\]

Impose

\[
\boxed{
T_j
\geq
C_j-d_j.
}
\]

Then minimizing

\[
\boxed{
\sum_jT_j
}
\]

forces

\[
T_j
=
\max
\{
C_j-d_j,
0
\}
\]

at optimality.

%%%%%%%%%%%%%%%%%%%%%%%%%%%%%%%%%%%%%%%%%%%%%%%%%%%%%%%%%%%%
\subsection{Time-Indexed Scheduling Formulation}
\label{subsec:time-indexed-scheduling}
%%%%%%%%%%%%%%%%%%%%%%%%%%%%%%%%%%%%%%%%%%%%%%%%%%%%%%%%%%%%

Suppose time is discretized.

Define

\[
x_{jt}
=
\begin{cases}
1,
&
\text{job }j\text{ starts at time }t,
\\
0,
&
\text{otherwise}.
\end{cases}
\]

Each job starts exactly once:

\[
\boxed{
\sum_t
x_{jt}
=
1.
}
\]

%%%%%%%%%%%%%%%%%%%%%%%%%%%%%%%%%%%%%%%%%%%%%%%%%%%%%%%%%%%%
\subsection{Time-Indexed Capacity Constraints}
\label{subsec:time-indexed-capacity}
%%%%%%%%%%%%%%%%%%%%%%%%%%%%%%%%%%%%%%%%%%%%%%%%%%%%%%%%%%%%

For one machine, at most one job may be active at time \(\tau\).

Thus,

\[
\boxed{
\sum_j
\sum_{
t:
t\leq\tau<t+p_j
}
x_{jt}
\leq
1.
}
\]

Time-indexed formulations can be strong but may require many binary
variables.

%%%%%%%%%%%%%%%%%%%%%%%%%%%%%%%%%%%%%%%%%%%%%%%%%%%%%%%%%%%%
\subsection{Release Dates in Time-Indexed Models}
\label{subsec:time-indexed-release-dates}
%%%%%%%%%%%%%%%%%%%%%%%%%%%%%%%%%%%%%%%%%%%%%%%%%%%%%%%%%%%%

If job \(j\) cannot start before \(r_j\), then simply omit variables with

\[
t<r_j,
\]

or impose

\[
\boxed{
x_{jt}=0,
\qquad
t<r_j.
}
\]

Similarly, a deadline can restrict allowable completion periods.

%%%%%%%%%%%%%%%%%%%%%%%%%%%%%%%%%%%%%%%%%%%%%%%%%%%%%%%%%%%%
\subsection{Sequence-Dependent Setup Times}
\label{subsec:sequence-dependent-setup-times}
%%%%%%%%%%%%%%%%%%%%%%%%%%%%%%%%%%%%%%%%%%%%%%%%%%%%%%%%%%%%

Suppose processing job \(j\) immediately after job \(i\) requires setup
time

\[
\boxed{
s_{ij}.
}
\]

Then if \(i\) precedes \(j\),

\[
\boxed{
S_j
\geq
C_i+s_{ij}.
}
\]

Examples include:

\[
\boxed{
\text{machine cleaning},
}
\]

\[
\boxed{
\text{tool changes},
}
\]

and

\[
\boxed{
\text{product changeovers}.
}
\]

%%%%%%%%%%%%%%%%%%%%%%%%%%%%%%%%%%%%%%%%%%%%%%%%%%%%%%%%%%%%
\subsection{Sequence-Independent Setup Times}
\label{subsec:sequence-independent-setup-times}
%%%%%%%%%%%%%%%%%%%%%%%%%%%%%%%%%%%%%%%%%%%%%%%%%%%%%%%%%%%%

If setup time depends only on the next job,

\[
s_j,
\]

then it may often be combined with processing time:

\[
\boxed{
p_j'
=
s_j+p_j.
}
\]

Sequence-dependent setup times cannot generally be absorbed this way.

%%%%%%%%%%%%%%%%%%%%%%%%%%%%%%%%%%%%%%%%%%%%%%%%%%%%%%%%%%%%
\subsection{Release Dates and Idle Time}
\label{subsec:release-dates-idle-time}
%%%%%%%%%%%%%%%%%%%%%%%%%%%%%%%%%%%%%%%%%%%%%%%%%%%%%%%%%%%%

Without release dates and with a regular objective, intentional idle time
on a single machine is usually unnecessary.

With release dates, idle time can be unavoidable.

In some problems, waiting for a future high-priority job can even be
strategically useful.

%%%%%%%%%%%%%%%%%%%%%%%%%%%%%%%%%%%%%%%%%%%%%%%%%%%%%%%%%%%%
\subsection{Shortest Remaining Processing Time}
\label{subsec:srpt}
%%%%%%%%%%%%%%%%%%%%%%%%%%%%%%%%%%%%%%%%%%%%%%%%%%%%%%%%%%%%

With preemption and release dates, the shortest remaining processing time
rule, abbreviated SRPT, always processes an available job with minimum
remaining processing time.

For the classical problem

\[
\boxed{
1
\mid
r_j,\mathrm{pmtn}
\mid
\sum C_j,
}
\]

SRPT is optimal.

%%%%%%%%%%%%%%%%%%%%%%%%%%%%%%%%%%%%%%%%%%%%%%%%%%%%%%%%%%%%
\subsection{Interpretation of SRPT}
\label{subsec:srpt-interpretation}
%%%%%%%%%%%%%%%%%%%%%%%%%%%%%%%%%%%%%%%%%%%%%%%%%%%%%%%%%%%%

When a shorter job arrives, it may interrupt a longer job.

Completing short jobs quickly reduces the number of unfinished jobs in
the system.

This is the preemptive analogue of SPT.

%%%%%%%%%%%%%%%%%%%%%%%%%%%%%%%%%%%%%%%%%%%%%%%%%%%%%%%%%%%%
\subsection{Scheduling With Deadlines}
\label{subsec:scheduling-deadlines}
%%%%%%%%%%%%%%%%%%%%%%%%%%%%%%%%%%%%%%%%%%%%%%%%%%%%%%%%%%%%

A deadline

\[
\bar d_j
\]

is a hard requirement:

\[
\boxed{
C_j
\leq
\bar d_j.
}
\]

A due date is usually a soft target that may be violated with a penalty.

Thus:

\[
\boxed{
\text{deadline}
\neq
\text{due date}.
}
\]

%%%%%%%%%%%%%%%%%%%%%%%%%%%%%%%%%%%%%%%%%%%%%%%%%%%%%%%%%%%%
\subsection{Feasibility Versus Optimization}
\label{subsec:scheduling-feasibility}
%%%%%%%%%%%%%%%%%%%%%%%%%%%%%%%%%%%%%%%%%%%%%%%%%%%%%%%%%%%%

Some scheduling problems ask only whether all timing requirements can be
satisfied.

Others seek the best feasible schedule.

Thus scheduling includes both:

\[
\boxed{
\text{feasibility problems}
}
\]

and

\[
\boxed{
\text{optimization problems}.
}
\]

%%%%%%%%%%%%%%%%%%%%%%%%%%%%%%%%%%%%%%%%%%%%%%%%%%%%%%%%%%%%
\subsection{Lower Bounds}
\label{subsec:scheduling-lower-bounds}
%%%%%%%%%%%%%%%%%%%%%%%%%%%%%%%%%%%%%%%%%%%%%%%%%%%%%%%%%%%%

Lower bounds are important in exact algorithms.

For makespan, useful bounds can arise from:

\[
\boxed{
\text{total workload},
}
\]

\[
\boxed{
\text{largest job},
}
\]

\[
\boxed{
\text{critical paths},
}
\]

or

\[
\boxed{
\text{machine workloads}.
}
\]

A lower bound helps determine how close a candidate schedule is to optimal.

%%%%%%%%%%%%%%%%%%%%%%%%%%%%%%%%%%%%%%%%%%%%%%%%%%%%%%%%%%%%
\subsection{Machine-Workload Bound}
\label{subsec:machine-workload-bound}
%%%%%%%%%%%%%%%%%%%%%%%%%%%%%%%%%%%%%%%%%%%%%%%%%%%%%%%%%%%%

Suppose every job shop operation assigned to machine \(i\) has processing
time \(p_{ij}\).

Any schedule must satisfy

\[
\boxed{
C_{\max}
\geq
\sum_j
p_{ij}.
}
\]

Therefore,

\[
\boxed{
C_{\max}
\geq
\max_i
\sum_jp_{ij}.
}
\]

%%%%%%%%%%%%%%%%%%%%%%%%%%%%%%%%%%%%%%%%%%%%%%%%%%%%%%%%%%%%
\subsection{Critical-Path Bound in Job Shops}
\label{subsec:job-shop-critical-path-bound}
%%%%%%%%%%%%%%%%%%%%%%%%%%%%%%%%%%%%%%%%%%%%%%%%%%%%%%%%%%%%

For each job, its operations must follow a specified route.

Therefore,

\[
\boxed{
C_{\max}
\geq
\sum_{k\in J_j}
p_{jk}
}
\]

for every job \(j\).

A lower bound is therefore

\[
\boxed{
\max_j
\{
\text{total processing requirement of job }j
\}.
}
\]

%%%%%%%%%%%%%%%%%%%%%%%%%%%%%%%%%%%%%%%%%%%%%%%%%%%%%%%%%%%%
\subsection{Branch-and-Bound for Scheduling}
\label{subsec:scheduling-branch-bound}
%%%%%%%%%%%%%%%%%%%%%%%%%%%%%%%%%%%%%%%%%%%%%%%%%%%%%%%%%%%%

Branch-and-bound can search over sequencing or assignment decisions.

A node may represent a partially fixed schedule.

At each node:

\begin{enumerate}
    \item compute a lower bound;
    \item construct or maintain a feasible incumbent;
    \item prune if the lower bound cannot improve the incumbent;
    \item otherwise branch on another scheduling decision.
\end{enumerate}

%%%%%%%%%%%%%%%%%%%%%%%%%%%%%%%%%%%%%%%%%%%%%%%%%%%%%%%%%%%%
\subsection{Dynamic Programming for Scheduling}
\label{subsec:scheduling-dynamic-programming}
%%%%%%%%%%%%%%%%%%%%%%%%%%%%%%%%%%%%%%%%%%%%%%%%%%%%%%%%%%%%

Dynamic programming can be effective when the state is manageable.

A subset state may represent:

\[
\boxed{
S
=
\text{set of jobs already completed}.
}
\]

The remaining schedule depends only on \(S\) and any additional timing
information required.

This can reduce repeated enumeration, although the state space may still
contain

\[
\boxed{
2^n
}
\]

subsets.

%%%%%%%%%%%%%%%%%%%%%%%%%%%%%%%%%%%%%%%%%%%%%%%%%%%%%%%%%%%%
\subsection{Local Search for Scheduling}
\label{subsec:scheduling-local-search}
%%%%%%%%%%%%%%%%%%%%%%%%%%%%%%%%%%%%%%%%%%%%%%%%%%%%%%%%%%%%

A scheduling neighborhood may use:

\[
\boxed{
\text{adjacent swaps},
}
\]

\[
\boxed{
\text{pairwise exchanges},
}
\]

\[
\boxed{
\text{job insertion},
}
\]

or

\[
\boxed{
\text{machine reassignment}.
}
\]

Local search repeatedly replaces the current schedule with a better nearby
schedule.

%%%%%%%%%%%%%%%%%%%%%%%%%%%%%%%%%%%%%%%%%%%%%%%%%%%%%%%%%%%%
\subsection{Adjacent Swap Neighborhood}
\label{subsec:adjacent-swap-neighborhood}
%%%%%%%%%%%%%%%%%%%%%%%%%%%%%%%%%%%%%%%%%%%%%%%%%%%%%%%%%%%%

For a sequence

\[
J_1,J_2,\ldots,J_n,
\]

an adjacent-swap move exchanges

\[
J_k
\]

and

\[
J_{k+1}.
\]

There are only

\[
\boxed{
n-1
}
\]

adjacent swaps, making this a small neighborhood.

%%%%%%%%%%%%%%%%%%%%%%%%%%%%%%%%%%%%%%%%%%%%%%%%%%%%%%%%%%%%
\subsection{Insertion Neighborhood}
\label{subsec:insertion-neighborhood}
%%%%%%%%%%%%%%%%%%%%%%%%%%%%%%%%%%%%%%%%%%%%%%%%%%%%%%%%%%%%

An insertion move removes one job from its current position and inserts it
elsewhere.

This neighborhood is larger than adjacent swaps and can escape some local
optima that adjacent swaps cannot.

%%%%%%%%%%%%%%%%%%%%%%%%%%%%%%%%%%%%%%%%%%%%%%%%%%%%%%%%%%%%
\subsection{Dispatching Versus Offline Scheduling}
\label{subsec:dispatching-versus-offline}
%%%%%%%%%%%%%%%%%%%%%%%%%%%%%%%%%%%%%%%%%%%%%%%%%%%%%%%%%%%%

Offline scheduling assumes the relevant jobs and information are known
before scheduling begins.

Online scheduling receives jobs over time.

A dispatching rule makes repeated local decisions as information becomes
available.

Examples include:

\[
\boxed{
\text{SPT},
}
\]

\[
\boxed{
\text{EDD},
}
\]

and

\[
\boxed{
\text{critical-ratio dispatching}.
}
\]

%%%%%%%%%%%%%%%%%%%%%%%%%%%%%%%%%%%%%%%%%%%%%%%%%%%%%%%%%%%%
\subsection{Online Scheduling}
\label{subsec:online-scheduling}
%%%%%%%%%%%%%%%%%%%%%%%%%%%%%%%%%%%%%%%%%%%%%%%%%%%%%%%%%%%%

In online scheduling, future jobs may be unknown.

An online algorithm must make decisions without full future information.

Performance can be evaluated using competitive analysis, which compares
the online schedule against an optimal schedule that knows the entire
input in advance.

%%%%%%%%%%%%%%%%%%%%%%%%%%%%%%%%%%%%%%%%%%%%%%%%%%%%%%%%%%%%
\subsection{Batch Scheduling}
\label{subsec:batch-scheduling}
%%%%%%%%%%%%%%%%%%%%%%%%%%%%%%%%%%%%%%%%%%%%%%%%%%%%%%%%%%%%

Some machines can process several compatible jobs simultaneously as one
batch.

A batch has:

\[
\boxed{
\text{capacity restrictions}
}
\]

and may have a processing time determined by the jobs assigned to it.

Batch scheduling occurs in ovens, furnaces, chemical processing, and data
processing.

%%%%%%%%%%%%%%%%%%%%%%%%%%%%%%%%%%%%%%%%%%%%%%%%%%%%%%%%%%%%
\subsection{No-Wait Scheduling}
\label{subsec:no-wait-scheduling}
%%%%%%%%%%%%%%%%%%%%%%%%%%%%%%%%%%%%%%%%%%%%%%%%%%%%%%%%%%%%

In a no-wait flow shop, once a job starts, it must continue through its
sequence of machines without waiting between operations.

This occurs when intermediate storage is impossible or product quality
would deteriorate.

No-wait constraints tightly couple operation start times.

%%%%%%%%%%%%%%%%%%%%%%%%%%%%%%%%%%%%%%%%%%%%%%%%%%%%%%%%%%%%
\subsection{Blocking Scheduling}
\label{subsec:blocking-scheduling}
%%%%%%%%%%%%%%%%%%%%%%%%%%%%%%%%%%%%%%%%%%%%%%%%%%%%%%%%%%%%

If there is no buffer between machines, a completed job may remain on its
current machine until the next machine becomes available.

This is called blocking.

The upstream machine then cannot begin another job.

%%%%%%%%%%%%%%%%%%%%%%%%%%%%%%%%%%%%%%%%%%%%%%%%%%%%%%%%%%%%
\subsection{Workforce Scheduling}
\label{subsec:workforce-scheduling}
%%%%%%%%%%%%%%%%%%%%%%%%%%%%%%%%%%%%%%%%%%%%%%%%%%%%%%%%%%%%

Workforce scheduling assigns employees to shifts or tasks while satisfying
coverage requirements.

Let

\[
x_s
=
\text{number of workers assigned to shift }s.
\]

If period \(t\) requires at least \(d_t\) workers,

\[
\boxed{
\sum_s
a_{ts}x_s
\geq
d_t,
}
\]

where \(a_{ts}=1\) if shift \(s\) covers period \(t\).

%%%%%%%%%%%%%%%%%%%%%%%%%%%%%%%%%%%%%%%%%%%%%%%%%%%%%%%%%%%%
\subsection{Shift Scheduling}
\label{subsec:shift-scheduling}
%%%%%%%%%%%%%%%%%%%%%%%%%%%%%%%%%%%%%%%%%%%%%%%%%%%%%%%%%%%%

A basic shift-scheduling model is

\[
\boxed{
\begin{aligned}
\min
\quad&
\sum_s
c_sx_s
\\
\text{subject to}
\quad&
\sum_s
a_{ts}x_s
\geq
d_t,
\qquad
t=1,\ldots,T,
\\
&
x_s\in\mathbb{Z}_{\geq0}.
\end{aligned}
}
\]

More detailed workforce models include:

\[
\boxed{
\text{skills},
}
\]

\[
\boxed{
\text{breaks},
}
\]

\[
\boxed{
\text{overtime},
}
\]

and

\[
\boxed{
\text{fairness}.
}
\]

%%%%%%%%%%%%%%%%%%%%%%%%%%%%%%%%%%%%%%%%%%%%%%%%%%%%%%%%%%%%
\subsection{Crew Scheduling}
\label{subsec:crew-scheduling}
%%%%%%%%%%%%%%%%%%%%%%%%%%%%%%%%%%%%%%%%%%%%%%%%%%%%%%%%%%%%

Crew scheduling selects feasible sequences of work activities.

Each complete feasible duty can be treated as a column.

This often leads to large set-covering or set-partitioning models.

Column generation is therefore central in airline, rail, and transit crew
scheduling.

%%%%%%%%%%%%%%%%%%%%%%%%%%%%%%%%%%%%%%%%%%%%%%%%%%%%%%%%%%%%
\subsection{Transportation Scheduling}
\label{subsec:transportation-scheduling}
%%%%%%%%%%%%%%%%%%%%%%%%%%%%%%%%%%%%%%%%%%%%%%%%%%%%%%%%%%%%

Transportation scheduling may involve:

\[
\boxed{
\text{vehicles},
}
\]

\[
\boxed{
\text{drivers},
}
\]

\[
\boxed{
\text{gates},
}
\]

\[
\boxed{
\text{runways},
}
\]

or

\[
\boxed{
\text{loading resources}.
}
\]

Many such problems combine:

\[
\boxed{
\text{network optimization}
+
\text{integer programming}
+
\text{scheduling}.
}
\]

%%%%%%%%%%%%%%%%%%%%%%%%%%%%%%%%%%%%%%%%%%%%%%%%%%%%%%%%%%%%
\subsection{Machine Availability}
\label{subsec:machine-availability}
%%%%%%%%%%%%%%%%%%%%%%%%%%%%%%%%%%%%%%%%%%%%%%%%%%%%%%%%%%%%

A machine may be unavailable during maintenance or other downtime.

If machine \(i\) is unavailable on interval

\[
[a,b],
\]

then no nonpreemptive operation assigned to that machine may overlap the
interval.

Calendar constraints can therefore become part of scheduling feasibility.

%%%%%%%%%%%%%%%%%%%%%%%%%%%%%%%%%%%%%%%%%%%%%%%%%%%%%%%%%%%%
\subsection{Maintenance Scheduling}
\label{subsec:maintenance-scheduling}
%%%%%%%%%%%%%%%%%%%%%%%%%%%%%%%%%%%%%%%%%%%%%%%%%%%%%%%%%%%%

Maintenance scheduling balances:

\[
\boxed{
\text{equipment availability},
}
\]

\[
\boxed{
\text{maintenance requirements},
}
\]

and

\[
\boxed{
\text{production demand}.
}
\]

Preventive maintenance can reduce failure risk but temporarily removes
capacity.

This creates a natural optimization tradeoff.

%%%%%%%%%%%%%%%%%%%%%%%%%%%%%%%%%%%%%%%%%%%%%%%%%%%%%%%%%%%%
\subsection{Stochastic Scheduling Preview}
\label{subsec:stochastic-scheduling-preview}
%%%%%%%%%%%%%%%%%%%%%%%%%%%%%%%%%%%%%%%%%%%%%%%%%%%%%%%%%%%%

Real processing times may be random.

Machines may fail.

Jobs may arrive unexpectedly.

Then the schedule must respond to uncertainty.

Approaches include:

\[
\boxed{
\text{stochastic optimization},
}
\]

\[
\boxed{
\text{robust scheduling},
}
\]

\[
\boxed{
\text{dynamic rescheduling},
}
\]

and

\[
\boxed{
\text{dispatching policies}.
}
\]

%%%%%%%%%%%%%%%%%%%%%%%%%%%%%%%%%%%%%%%%%%%%%%%%%%%%%%%%%%%%
\subsection{Robust Scheduling Preview}
\label{subsec:robust-scheduling-preview}
%%%%%%%%%%%%%%%%%%%%%%%%%%%%%%%%%%%%%%%%%%%%%%%%%%%%%%%%%%%%

A robust schedule attempts to remain effective under uncertain processing
times or disruptions.

One may include:

\[
\boxed{
\text{time buffers},
}
\]

\[
\boxed{
\text{slack},
}
\]

or

\[
\boxed{
\text{uncertainty sets}.
}
\]

The goal is often to trade nominal efficiency for resilience.

%%%%%%%%%%%%%%%%%%%%%%%%%%%%%%%%%%%%%%%%%%%%%%%%%%%%%%%%%%%%
\subsection{Schedule Stability}
\label{subsec:schedule-stability}
%%%%%%%%%%%%%%%%%%%%%%%%%%%%%%%%%%%%%%%%%%%%%%%%%%%%%%%%%%%%

When disruptions occur, repeatedly replacing the original schedule may be
operationally expensive.

A rescheduling objective may therefore penalize deviation from the original
plan:

\[
\boxed{
\text{new performance cost}
+
\lambda
(
\text{schedule-change cost}
).
}
\]

This distinguishes optimization quality from schedule stability.

%%%%%%%%%%%%%%%%%%%%%%%%%%%%%%%%%%%%%%%%%%%%%%%%%%%%%%%%%%%%
\subsection{Scheduling Complexity}
\label{subsec:scheduling-complexity}
%%%%%%%%%%%%%%%%%%%%%%%%%%%%%%%%%%%%%%%%%%%%%%%%%%%%%%%%%%%%

Scheduling problems range from very easy to computationally difficult.

Examples with simple optimal rules include:

\[
\boxed{
1\mid\mid\sum C_j
}
\]

using SPT,

and

\[
\boxed{
1\mid\mid L_{\max}
}
\]

using EDD.

Other problems, such as general job-shop makespan scheduling, are NP-hard.

%%%%%%%%%%%%%%%%%%%%%%%%%%%%%%%%%%%%%%%%%%%%%%%%%%%%%%%%%%%%
\subsection{Why Small Model Changes Matter}
\label{subsec:scheduling-small-model-changes}
%%%%%%%%%%%%%%%%%%%%%%%%%%%%%%%%%%%%%%%%%%%%%%%%%%%%%%%%%%%%

Adding only one feature can dramatically change complexity.

For example:

\[
\boxed{
\text{one machine}
}
\]

may be easy under one objective,

while adding:

\[
\boxed{
\text{release dates},
}
\]

\[
\boxed{
\text{sequence-dependent setups},
}
\]

or

\[
\boxed{
\text{multiple machines}
}
\]

can create a much harder problem.

%%%%%%%%%%%%%%%%%%%%%%%%%%%%%%%%%%%%%%%%%%%%%%%%%%%%%%%%%%%%
\subsection{Scheduling Optimization Workflow}
\label{subsec:scheduling-workflow}
%%%%%%%%%%%%%%%%%%%%%%%%%%%%%%%%%%%%%%%%%%%%%%%%%%%%%%%%%%%%

A practical workflow is:

\begin{enumerate}
    \item identify jobs and operations;
    \item identify machines and resources;
    \item record processing times;
    \item record release dates and due dates;
    \item identify precedence requirements;
    \item determine whether preemption is allowed;
    \item identify setup times;
    \item identify machine availability restrictions;
    \item define the objective;
    \item write the problem using three-field notation when possible;
    \item check for a known optimal sequencing rule;
    \item derive lower bounds;
    \item construct a feasible schedule;
    \item choose exact optimization, approximation, or heuristic methods;
    \item compute completion-time performance measures;
    \item compare against lower bounds or optimality gaps;
    \item perform sensitivity analysis;
    \item consider robustness and rescheduling needs.
\end{enumerate}

%%%%%%%%%%%%%%%%%%%%%%%%%%%%%%%%%%%%%%%%%%%%%%%%%%%%%%%%%%%%
\subsection{Choosing a Scheduling Method}
\label{subsec:choosing-scheduling-method}
%%%%%%%%%%%%%%%%%%%%%%%%%%%%%%%%%%%%%%%%%%%%%%%%%%%%%%%%%%%%

For

\[
1\mid\mid\sum C_j,
\]

use:

\[
\boxed{
\text{SPT}.
}
\]

For

\[
1\mid\mid\sum w_jC_j,
\]

use:

\[
\boxed{
\text{Smith's rule}.
}
\]

For

\[
1\mid\mid L_{\max},
\]

use:

\[
\boxed{
\text{EDD}.
}
\]

For

\[
1\mid\mid\sum U_j,
\]

use:

\[
\boxed{
\text{Moore--Hodgson}.
}
\]

For

\[
F_2\mid\mid C_{\max},
\]

use:

\[
\boxed{
\text{Johnson's rule}.
}
\]

For difficult general scheduling problems, consider:

\[
\boxed{
\text{MILP},
}
\]

\[
\boxed{
\text{branch-and-bound},
}
\]

\[
\boxed{
\text{constraint programming},
}
\]

or

\[
\boxed{
\text{metaheuristics}.
}
\]

%%%%%%%%%%%%%%%%%%%%%%%%%%%%%%%%%%%%%%%%%%%%%%%%%%%%%%%%%%%%
\subsection{Common Errors}
\label{subsec:scheduling-common-errors}
%%%%%%%%%%%%%%%%%%%%%%%%%%%%%%%%%%%%%%%%%%%%%%%%%%%%%%%%%%%%

\subsubsection{Confusing Lateness and Tardiness}

Lateness can be negative:

\[
\boxed{
L_j=C_j-d_j.
}
\]

Tardiness cannot:

\[
\boxed{
T_j=\max\{L_j,0\}.
}
\]

%%%%%%%%%%%%%%%%%%%%%%%%%%%%%%%%%%%%%%%%%%%%%%%%%%%%%%%%%%%%
\subsubsection{Confusing Due Dates and Deadlines}

A due date is generally a target.

A deadline is generally a hard feasibility requirement.

They should not be modeled identically unless the application requires it.

%%%%%%%%%%%%%%%%%%%%%%%%%%%%%%%%%%%%%%%%%%%%%%%%%%%%%%%%%%%%
\subsubsection{Using SPT for Every Objective}

SPT minimizes total completion time in the classical single-machine
problem.

It does not generally minimize:

\[
\boxed{
L_{\max},
}
\]

\[
\boxed{
\sum T_j,
}
\]

or

\[
\boxed{
C_{\max}
}
\]

in unrelated scheduling environments.

%%%%%%%%%%%%%%%%%%%%%%%%%%%%%%%%%%%%%%%%%%%%%%%%%%%%%%%%%%%%
\subsubsection{Using EDD to Minimize Total Tardiness}

EDD minimizes maximum lateness for the classical single-machine problem.

It does not generally minimize total tardiness.

%%%%%%%%%%%%%%%%%%%%%%%%%%%%%%%%%%%%%%%%%%%%%%%%%%%%%%%%%%%%
\subsubsection{Ignoring Release Dates}

A schedule beginning job \(j\) before

\[
r_j
\]

is infeasible.

%%%%%%%%%%%%%%%%%%%%%%%%%%%%%%%%%%%%%%%%%%%%%%%%%%%%%%%%%%%%
\subsubsection{Ignoring Setup Times}

If changing from job \(i\) to job \(j\) requires setup time

\[
s_{ij},
\]

then omitting it can underestimate completion times substantially.

%%%%%%%%%%%%%%%%%%%%%%%%%%%%%%%%%%%%%%%%%%%%%%%%%%%%%%%%%%%%
\subsubsection{Allowing Machine Overlap}

Two nonpreemptive operations requiring the same single-capacity machine
cannot overlap.

This must be enforced explicitly.

%%%%%%%%%%%%%%%%%%%%%%%%%%%%%%%%%%%%%%%%%%%%%%%%%%%%%%%%%%%%
\subsubsection{Allowing a Job to Use Two Machines Simultaneously}

In ordinary shop scheduling, a job normally processes one operation at a
time unless the technological model explicitly allows simultaneous
operations.

%%%%%%%%%%%%%%%%%%%%%%%%%%%%%%%%%%%%%%%%%%%%%%%%%%%%%%%%%%%%
\subsubsection{Ignoring Precedence Constraints}

A schedule can satisfy all machine capacities and still be infeasible if
job routing or project precedence is violated.

%%%%%%%%%%%%%%%%%%%%%%%%%%%%%%%%%%%%%%%%%%%%%%%%%%%%%%%%%%%%
\subsubsection{Using an Excessively Large Big-\(M\)}

Very large \(M\) values weaken mixed-integer formulations and can damage
numerical performance.

Use the tightest justified scheduling horizon.

%%%%%%%%%%%%%%%%%%%%%%%%%%%%%%%%%%%%%%%%%%%%%%%%%%%%%%%%%%%%
\subsubsection{Assuming a Lower Bound Is a Feasible Schedule}

A lower bound estimates what no schedule can beat.

It need not itself be attainable.

%%%%%%%%%%%%%%%%%%%%%%%%%%%%%%%%%%%%%%%%%%%%%%%%%%%%%%%%%%%%
\subsubsection{Assuming a Heuristic Solution Is Optimal}

Dispatching rules, local search, and metaheuristics can produce good
schedules without proving optimality.

A valid lower bound or exact algorithm is required for a certificate.

%%%%%%%%%%%%%%%%%%%%%%%%%%%%%%%%%%%%%%%%%%%%%%%%%%%%%%%%%%%%
\subsubsection{Assuming Critical Path Alone Solves Resource-Constrained Projects}

CPM handles precedence timing.

If several activities compete for scarce resources, additional scheduling
decisions are required.

%%%%%%%%%%%%%%%%%%%%%%%%%%%%%%%%%%%%%%%%%%%%%%%%%%%%%%%%%%%%
\subsubsection{Interpreting PERT Estimates as Exact Distributions}

Classical PERT expected-time and variance formulas are approximations based
on modeling assumptions.

Real project durations may not follow those assumptions.

%%%%%%%%%%%%%%%%%%%%%%%%%%%%%%%%%%%%%%%%%%%%%%%%%%%%%%%%%%%%
\subsubsection{Ignoring Rescheduling Cost}

A mathematically improved schedule may be operationally undesirable if it
requires large last-minute changes.

Real systems may require schedule-stability objectives.

%%%%%%%%%%%%%%%%%%%%%%%%%%%%%%%%%%%%%%%%%%%%%%%%%%%%%%%%%%%%
\subsection{Applications}
\label{subsec:scheduling-applications}
%%%%%%%%%%%%%%%%%%%%%%%%%%%%%%%%%%%%%%%%%%%%%%%%%%%%%%%%%%%%

\begin{applicationbox}

\textbf{Manufacturing.}
Jobs must be sequenced across machines while respecting processing routes,
setups, due dates, and capacity.

\medskip

\textbf{Air Transportation.}
Aircraft, gates, crews, loading equipment, maintenance resources, and
runways all require coordinated schedules.

\medskip

\textbf{Computing.}
Processors schedule tasks to reduce completion time, response time,
lateness, or energy use.

\medskip

\textbf{Healthcare.}
Operating rooms, physicians, nurses, diagnostic equipment, and patient
appointments require scheduling under limited capacity.

\medskip

\textbf{Project Management.}
CPM, PERT, and resource-constrained scheduling coordinate activities with
precedence relationships.

\medskip

\textbf{Workforce Planning.}
Employees are assigned to shifts while satisfying coverage, qualification,
labor, and fairness requirements.

\medskip

\textbf{Logistics.}
Vehicles, docks, loading bays, warehouse resources, and delivery windows
create complex scheduling problems.

\medskip

\textbf{Maintenance.}
Preventive and corrective maintenance must be scheduled while preserving
operational availability.

\end{applicationbox}

%%%%%%%%%%%%%%%%%%%%%%%%%%%%%%%%%%%%%%%%%%%%%%%%%%%%%%%%%%%%
\subsection{Exercises}
\label{subsec:scheduling-exercises}
%%%%%%%%%%%%%%%%%%%%%%%%%%%%%%%%%%%%%%%%%%%%%%%%%%%%%%%%%%%%

\begin{exercise}[1]
Define processing time.
\end{exercise}

\begin{exercise}[1]
Define release date.
\end{exercise}

\begin{exercise}[1]
Define due date.
\end{exercise}

\begin{exercise}[1]
Define completion time.
\end{exercise}

\begin{exercise}[1]
Define flow time.
\end{exercise}

\begin{exercise}[1]
Define lateness.
\end{exercise}

\begin{exercise}[1]
Define tardiness.
\end{exercise}

\begin{exercise}[1]
Define makespan.
\end{exercise}

\begin{exercise}[1]
Define preemption.
\end{exercise}

\begin{exercise}[1]
Interpret

\[
1\mid\mid\sum C_j.
\]
\end{exercise}

\begin{exercise}[2]
For

\[
C_j=14,
\qquad
d_j=10,
\]

compute lateness and tardiness.
\end{exercise}

\begin{exercise}[2]
For

\[
C_j=8,
\qquad
d_j=12,
\]

compute lateness, tardiness, and earliness.
\end{exercise}

\begin{exercise}[2]
For

\[
r_j=3,
\qquad
C_j=11,
\]

compute flow time.
\end{exercise}

\begin{exercise}[2]
Order processing times

\[
7,2,9,4
\]

according to SPT.
\end{exercise}

\begin{exercise}[2]
Order due dates

\[
12,5,9,18
\]

according to EDD.
\end{exercise}

\begin{exercise}[2]
Compute

\[
\rho_j
=
\frac{p_j}{w_j}
\]

for several jobs and determine Smith's-rule order.
\end{exercise}

\begin{exercise}[3]
Use SPT to minimize total completion time for five jobs.
\end{exercise}

\begin{exercise}[3]
Compute every \(C_j\) and verify the resulting value of

\[
\sum C_j.
\]
\end{exercise}

\begin{exercise}[3]
Use Smith's rule to minimize

\[
\sum w_jC_j
\]

for a small job set.
\end{exercise}

\begin{exercise}[3]
Use EDD to minimize maximum lateness.
\end{exercise}

\begin{exercise}[3]
Compute \(L_j\), \(T_j\), and \(L_{\max}\) for the resulting schedule.
\end{exercise}

\begin{exercise}[3]
Apply the Moore--Hodgson algorithm to a small set of jobs.
\end{exercise}

\begin{exercise}[3]
Determine the minimum number of tardy jobs in the previous problem.
\end{exercise}

\begin{exercise}[3]
Construct a precedence graph from a list of precedence relations.
\end{exercise}

\begin{exercise}[3]
Perform a forward pass through a project network.
\end{exercise}

\begin{exercise}[3]
Perform a backward pass and compute activity slack.
\end{exercise}

\begin{exercise}[3]
Identify the critical path.
\end{exercise}

\begin{exercise}[3]
Compute a PERT expected duration from optimistic, most-likely, and
pessimistic values.
\end{exercise}

\begin{exercise}[3]
Compute the corresponding PERT activity variance.
\end{exercise}

\begin{exercise}[3]
For three identical parallel machines, compute the workload lower bound
for makespan.
\end{exercise}

\begin{exercise}[3]
Apply arbitrary list scheduling to a set of jobs.
\end{exercise}

\begin{exercise}[3]
Apply LPT scheduling to the same job set.
\end{exercise}

\begin{exercise}[3]
Compare the two makespans.
\end{exercise}

\begin{exercise}[3]
Determine whether the lower bound proves the LPT schedule optimal.
\end{exercise}

\begin{exercise}[3]
Apply Johnson's rule to a four-job, two-machine flow shop.
\end{exercise}

\begin{exercise}[3]
Construct the machine completion-time table for the Johnson schedule.
\end{exercise}

\begin{exercise}[3]
Compute the resulting flow-shop makespan.
\end{exercise}

\begin{exercise}[3]
Write disjunctive constraints for two jobs requiring the same machine.
\end{exercise}

\begin{exercise}[3]
Convert the previous disjunction to a Big-\(M\) formulation.
\end{exercise}

\begin{exercise}[3]
Introduce a makespan variable into a scheduling MILP.
\end{exercise}

\begin{exercise}[3]
Linearize

\[
T_j
=
\max
\{
C_j-d_j,0
\}.
\]
\end{exercise}

\begin{exercise}[3]
Construct a time-indexed formulation for three jobs.
\end{exercise}

\begin{exercise}[3]
Write one-machine capacity constraints in time-indexed form.
\end{exercise}

\begin{exercise}[3]
Add release dates to the previous formulation.
\end{exercise}

\begin{exercise}[3]
Model sequence-dependent setup times.
\end{exercise}

\begin{exercise}[3]
Compute slack and critical ratio for a job at a specified current time.
\end{exercise}

\begin{exercise}[3]
Explain when intentional idle time may be useful.
\end{exercise}

\begin{exercise}[3]
Explain the difference between offline and online scheduling.
\end{exercise}

\begin{exercise}[3]
Construct a basic workforce shift-scheduling model.
\end{exercise}

\begin{exercise}[3]
Construct a resource-capacity constraint for a project schedule.
\end{exercise}

\begin{exercise}[3]
Explain why RCPSP cannot generally be solved by CPM alone.
\end{exercise}

\begin{exercise}[3]
Construct lower bounds for a small job-shop makespan problem.
\end{exercise}

\begin{exercise}[3]
Explain how branch-and-bound can use these lower bounds.
\end{exercise}

\begin{exercise}[4]
Prove the SPT optimality theorem using adjacent interchange.
\end{exercise}

\begin{exercise}[4]
Prove Smith's rule using pairwise interchange.
\end{exercise}

\begin{exercise}[4]
Prove EDD optimality for

\[
1\mid\mid L_{\max}.
\]
\end{exercise}

\begin{exercise}[4]
Prove correctness of the Moore--Hodgson algorithm.
\end{exercise}

\begin{exercise}[4]
Prove the list-scheduling bound

\[
C_{\max}^{LS}
\leq
\left(
2-\frac1m
\right)
C_{\max}^*.
\]
\end{exercise}

\begin{exercise}[4]
Study the LPT approximation bound.
\end{exercise}

\begin{exercise}[4]
Prove Johnson's rule for

\[
F_2\mid\mid C_{\max}.
\]
\end{exercise}

\begin{exercise}[4]
Study single-machine total-tardiness dynamic programming.
\end{exercise}

\begin{exercise}[4]
Study

\[
1\mid r_j,\mathrm{pmtn}\mid\sum C_j
\]

and prove optimality of SRPT.
\end{exercise}

\begin{exercise}[4]
Study scheduling with precedence constraints.
\end{exercise}

\begin{exercise}[4]
Investigate Lawler's algorithm for precedence-constrained maximum lateness.
\end{exercise}

\begin{exercise}[4]
Study exact branch-and-bound methods for job-shop scheduling.
\end{exercise}

\begin{exercise}[4]
Study disjunctive-graph formulations of job shops.
\end{exercise}

\begin{exercise}[4]
Investigate time-indexed formulations and LP relaxation strength.
\end{exercise}

\begin{exercise}[4]
Study constraint programming for scheduling.
\end{exercise}

\begin{exercise}[4]
Study interval variables and global no-overlap constraints.
\end{exercise}

\begin{exercise}[4]
Investigate local-search neighborhoods for job-shop scheduling.
\end{exercise}

\begin{exercise}[4]
Study tabu search for scheduling.
\end{exercise}

\begin{exercise}[4]
Study simulated annealing for scheduling.
\end{exercise}

\begin{exercise}[4]
Investigate crew scheduling using set partitioning and column generation.
\end{exercise}

\begin{exercise}[4]
Study the Resource-Constrained Project Scheduling Problem.
\end{exercise}

\begin{exercise}[4]
Study robust project scheduling.
\end{exercise}

\begin{exercise}[4]
Study stochastic scheduling with random processing times.
\end{exercise}

\begin{exercise}[4]
Investigate competitive analysis for online scheduling.
\end{exercise}

%%%%%%%%%%%%%%%%%%%%%%%%%%%%%%%%%%%%%%%%%%%%%%%%%%%%%%%%%%%%
\subsection{Conceptual Summary}
\label{subsec:scheduling-summary}
%%%%%%%%%%%%%%%%%%%%%%%%%%%%%%%%%%%%%%%%%%%%%%%%%%%%%%%%%%%%

Scheduling determines how jobs use limited resources over time.

Important job parameters include:

\[
\boxed{
p_j,
\quad
r_j,
\quad
d_j.
}
\]

Important performance measures include:

\[
\boxed{
C_j,
}
\]

\[
\boxed{
F_j=C_j-r_j,
}
\]

\[
\boxed{
L_j=C_j-d_j,
}
\]

\[
\boxed{
T_j=\max\{C_j-d_j,0\},
}
\]

and

\[
\boxed{
C_{\max}=\max_jC_j.
}
\]

For classical single-machine problems:

\[
\boxed{
1\mid\mid\sum C_j
\rightarrow
\text{SPT},
}
\]

\[
\boxed{
1\mid\mid\sum w_jC_j
\rightarrow
\text{Smith's rule},
}
\]

\[
\boxed{
1\mid\mid L_{\max}
\rightarrow
\text{EDD},
}
\]

and

\[
\boxed{
1\mid\mid\sum U_j
\rightarrow
\text{Moore--Hodgson}.
}
\]

For two-machine flow-shop makespan:

\[
\boxed{
F_2\mid\mid C_{\max}
\rightarrow
\text{Johnson's rule}.
}
\]

For parallel machines, basic makespan bounds include

\[
\boxed{
C_{\max}
\geq
\max
\left\{
\frac{\sum_jp_j}{m},
\max_jp_j
\right\}.
}
\]

Scheduling theory therefore combines

\[
\boxed{
\text{optimization}
+
\text{combinatorics}
+
\text{algorithms}
+
\text{resource allocation}
+
\text{time}.
}
\]

%%%%%%%%%%%%%%%%%%%%%%%%%%%%%%%%%%%%%%%%%%%%%%%%%%%%%%%%%%%%
\subsection{Section Connections}
\label{subsec:scheduling-connections}
%%%%%%%%%%%%%%%%%%%%%%%%%%%%%%%%%%%%%%%%%%%%%%%%%%%%%%%%%%%%

\chapterconnection{
Scheduling Theory combines the discrete decision structures of
Section~\ref{sec:combinatorial-optimization}, the integer formulations of
Section~\ref{sec:integer-programming}, and the congestion perspective of
Section~\ref{sec:queueing-theory}. Classical rules such as SPT, EDD,
Smith's rule, Moore--Hodgson, LPT, and Johnson's rule exploit special
problem structure, while general shop and project scheduling often require
mixed-integer optimization, dynamic programming, branch-and-bound,
constraint programming, or heuristics. The next section, Inventory Theory,
studies another major operations-research problem: how much and when to
order when demand, holding cost, shortage cost, and replenishment interact
over time.
}

%%%%%%%%%%%%%%%%%%%%%%%%%%%%%%%%%%%%%%%%%%%%%%%%%%%%%%%%%%%%
% SECTION SOURCE TRACKING
%%%%%%%%%%%%%%%%%%%%%%%%%%%%%%%%%%%%%%%%%%%%%%%%%%%%%%%%%%%%

%------------------------------------------------------------
% 10.12 Scheduling Theory
%------------------------------------------------------------
%
% Source basis:
%   - Standard deterministic scheduling theory.
%   - Standard project-scheduling and operations-research methods.
%   - Standard machine-scheduling algorithms and formulations.
%
% Used for:
%   - jobs, machines, and processing times;
%   - release dates and due dates;
%   - start and completion times;
%   - flow time;
%   - lateness, tardiness, and earliness;
%   - tardy-job indicators;
%   - makespan;
%   - weighted completion time;
%   - regular and nonregular objectives;
%   - preemption;
%   - single, parallel, flow-shop, job-shop, and open-shop environments;
%   - Graham three-field notation;
%   - SPT;
%   - Smith's rule / WSPT;
%   - EDD;
%   - Moore--Hodgson;
%   - dispatching rules;
%   - slack and critical ratio;
%   - precedence constraints;
%   - precedence graphs;
%   - critical paths;
%   - CPM;
%   - PERT;
%   - RCPSP;
%   - parallel-machine makespan bounds;
%   - list scheduling;
%   - LPT;
%   - Johnson's rule;
%   - permutation flow shops;
%   - disjunctive constraints;
%   - Big-M sequencing models;
%   - makespan and tardiness linearizations;
%   - time-indexed formulations;
%   - sequence-dependent setups;
%   - SRPT;
%   - deadlines;
%   - branch-and-bound;
%   - dynamic programming;
%   - local search;
%   - online scheduling;
%   - batch, blocking, and no-wait scheduling;
%   - workforce and crew scheduling;
%   - transportation scheduling;
%   - maintenance scheduling;
%   - stochastic and robust scheduling previews;
%   - schedule stability;
%   - applications and exercises.
%
% Notes:
%   - Inventory Theory is developed next in Section 10.13.
%   - Decision Theory follows in Section 10.14.
%   - Detailed stochastic optimization appears in Section 10.8.
%
%%%%%%%%%%%%%%%%%%%%%%%%%%%%%%%%%%%%%%%%%%%%%%%%%%%%%%%%%%%%